% =====================================================================
\chapter*{Final Build Statement (v7.2)}
\addcontentsline{toc}{chapter}{Final Build Statement (v7.2)}

\noindent
This volume represents the \textbf{Integrated System v7.2} of the NRMO
research programme, designated as the architectural completion of the
defensive side of NRMO. It is published as a self-contained reference
covering the theory, implementation, and empirical validation of NRMO's
governance--execution architecture across both individual and collective
scales.

\section*{What v7.2 contains}

\begin{itemize}
\item \textbf{NRMO\_Origin}: the foundational veto-only formulation
  (4 fixed rules) preserved without modification (Part~I).
\item \textbf{NRMO}: the strengthened operational form with 13 enhancements
  A--M, Adaptive Tuning Layer, and $\Omega$ Full Strong Engine (Parts~VII--IX
  and~XIII, v6.4 chapter).
\item \textbf{NRMO\_Collective}: the sibling theory introducing collective
  governance with 6 P--U mechanisms (v7.0 chapter) and the
  Collective Strong Engine with 6 W--AB mechanisms plus 4-rule Collective
  Core veto (v7.2 chapter). This restores architectural symmetry: at
  every scale (individual / family / lineage / civilisation), both veto
  and search components coexist.
\item \textbf{World Simulation suite}: a 9-civilisation, 75-generation,
  3000-year multi-agent simulator with cohort spotlighting, family
  lineage chronicling, inter-civilisation interaction, memetic dynamics,
  black-swan events, counterfactual reasoning, and calibration to
  historical demography. The full v6.1 unified codebase implements
  all three NRMO controllers.
\item \textbf{v52 reference codebase}: the original single-civilisation
  research implementation from which v6.3 onward was bridged.
\end{itemize}

\section*{What v7.2 deliberately does not contain}

\subsection*{Vision (North Star concept)}

A specification proposed during v7.2 development would have added a
``Common Vision'' layer to both individual and collective Strong
Engines, formalising the phrase
``\textit{Walk and observe the world without forgetting; look up to the North Star.}''
Implementation was rejected by the architectural reasoning recorded
during the design discussion. The decisive insight:

\begin{quote}
\itshape
NRMO must not possess Vision.\\
Vision belongs to the user of NRMO, not to NRMO itself.\\
NRMO is a decision aid framework, not a value system.
\end{quote}

This insight repositions NRMO definitively as a \emph{decision aid}
(neutral to the user's Vision) rather than a \emph{value system}
(prescribing direction). The Authority Hierarchy is therefore:

\[
\boxed{\text{Human Sovereign} \to \text{Vision (held by Human)} \to
       \text{NRMO} \to \text{Engines}}
\]

The user retains Vision; NRMO operates beneath Vision as admissibility
and ruin-rejection layer. This structural placement protects NRMO from
ideological drift and preserves the Spec's prohibition against
NRMO becoming ``doctrine, ideology, or replacement for human
sovereignty''.

\subsection*{Offensive theory (provisionally named NPMO)}

A complementary theory for the offensive side
(``Non-Passivity Maximizing Objective''), addressing active pursuit
rather than ruin-rejection, is recognised as a necessary counterpart
to NRMO but is reserved for future work. The intended architecture
positions NRMO and NPMO as siblings under a unifying Holdings layer:

\[
\boxed{\text{Holdings (user-level integration)} \to
       \text{NRMO (defensive) + NPMO (offensive)} \to \text{Engines}}
\]

NPMO's specification, implementation, and empirical validation are
out of scope for v7.2 and will appear in a separate volume.

\section*{Empirical position}

Three-seed averaged comparison (scale=small, 9 civilisations,
75 generations) summary, NRMO\_vNext strategy score:

\begin{center}
\small
\begin{tabular}{lrrrrr}
\toprule
World & Baseline (v6.0) & NRMO (v6.4) & NRMO\_Coll. (v7.0) & NRMO\_Coll. (v7.1) & Faith\_Communal \\
\midrule
Japan              & 2.148 & \textbf{2.170} & 1.729 & 1.699 & 1.997 \\
China              & 2.109 & \textbf{2.182} & 1.913 & 1.907 & 2.146 \\
Europe             & 0.769 & 0.835 & 1.108 & \textbf{1.118} & 1.450 \\
Islamic            & 1.544 & \textbf{1.616} & 1.395 & 1.372 & 1.431 \\
Indic              & 2.017 & \textbf{2.074} & 1.532 & 1.555 & 1.770 \\
Polynesian         & 0.943 & 1.012 & 1.156 & \textbf{1.207} & 1.585 \\
IndigenousAmericas & 0.213 & 0.195 & 0.837 & \textbf{0.890} & 0.752 \\
\bottomrule
\end{tabular}
\end{center}

\begin{itemize}
\item NRMO (v6.4) dominates in peaceful worlds.
\item NRMO\_Collective (v7.0/v7.1) dominates in catastrophic worlds.
\item In IndigenousAmericas, NRMO\_Collective exceeds Faith\_Communal
  ($0.890 > 0.752$): the first empirical case where a non-religious
  framework outperforms Faith-based collective survival.
\item The trade-off between peacetime efficiency and crisis-time
  survival is irreducible and operationally explicit.
\end{itemize}

\section*{Architectural invariants}

Throughout all extensions A--M, P--U, W--AB, the following invariants
hold without exception:

\begin{enumerate}
\item NRMO Core defines admissibility $A_t$; Strong Engines search
  within $A_t$. Neither engine may redefine NRMO boundaries.
\item Tuning Layer adjusts thresholds but never enters NRMO scoring.
\item Collective layer (P--U, W--AB) operates downstream of NRMO and
  $\Omega$ Full, never modifying veto logic.
\item Governance--execution separation is scale-invariant: applied
  identically at individual and collective levels.
\item NRMO does not possess Vision; the user does.
\end{enumerate}

\section*{This volume's intended use}

This volume is published as a self-contained reference for:
\begin{itemize}
\item Researchers studying decision architectures under absorbing
  failure risk;
\item Practitioners requiring veto + search separation in governance
  systems;
\item Future readers tracing the development history of NRMO theory.
\end{itemize}

It is \textbf{not} intended as:
\begin{itemize}
\item A philosophical doctrine to be adopted;
\item An ideology or value system;
\item A replacement for human judgement at any scale.
\end{itemize}

The reader is the sovereign. NRMO is the tool.
